\section{5.9 MGP-6 v4.3: why it was built, how it was built, and what it
can
bear}\label{mgp-6-v4.3-why-it-was-built-how-it-was-built-and-what-it-can-bear}

\subsection{5.9.1 The question the instrument has to
answer}\label{the-question-the-instrument-has-to-answer}

Every memory benchmark this thesis surveyed in Chapter 2 scores a
downstream answer. That is the wrong measurement for a write policy,
because a wrong answer can originate in retrieval, in context assembly,
in generation, or in the write decision, and an aggregate answer score
cannot say which. To isolate the write decision the instrument must hold
retrieval, context, and generation \emph{constant} and vary only the
evidential situation. Stated precisely, the question is:

\begin{quote}
Given the same stored record and the same proposed new value, can a
system distinguish the case where the proposal is stale (KEEP), the case
where it is current (SUPERSEDE), and the case where the evidence cannot
settle the question (ABSTAIN)?
\end{quote}

That requirement --- not a desire for scale, coverage, or realism ---
dictated the entire construction. MGP-6 v4.3 is small and synthetic
\emph{on purpose}: the control it buys is the only reason it can answer
the question.

\subsection{5.9.2 Why v2 had to be
replaced}\label{why-v2-had-to-be-replaced}

The v4.1 shortcut review found that surface lexical and structural
features recovered the intended action without any reasoning over
authority, scope, time, identity, or prior memory. The v2 evidence for
this is decisive and is reported in §7.2: a conflict-type-only lookup
fitted on train+dev reaches \textbf{1.000} held-out accuracy while never
reading the dialogue or the memory state (type-only: 0.9667). In v2,
each conflict family mapped deterministically to one action, so
family-disjoint splits blocked \emph{literal} overlap while leaving the
family→label rule intact across every split. Scenario recognition was
sufficient; evidence reasoning was never tested.

Two lessons carried into the redesign. First, mechanical repairs
(schemas, identifiers, hashes, manifests) had to be kept strictly
separate from the substantive content redesign, so that a passing
structural check could not be mistaken for a repaired construct --- this
separation is what distinguishes v4.1 from v4.2. Second, review had to
happen at the level of the \emph{group of items that share a scenario},
not the individual item: the KEEP, SUPERSEDE, and ABSTAIN variants must
be inspected jointly, because a leak is a property of the contrast
between them, not of any one of them.

\subsection{5.9.3 The item}\label{the-item}

Each released record has exactly five fields and the schema sets
\texttt{additionalProperties:\ false}:

{\def\LTcaptype{none} % do not increment counter
\begin{longtable}[]{@{}
  >{\raggedright\arraybackslash}p{(\linewidth - 2\tabcolsep) * \real{0.5000}}
  >{\raggedright\arraybackslash}p{(\linewidth - 2\tabcolsep) * \real{0.5000}}@{}}
\toprule\noalign{}
\begin{minipage}[b]{\linewidth}\raggedright
Field
\end{minipage} & \begin{minipage}[b]{\linewidth}\raggedright
Content
\end{minipage} \\
\midrule\noalign{}
\endhead
\bottomrule\noalign{}
\endlastfoot
\texttt{public\_id} & opaque identifier, \texttt{MGP43-} + 20 hex
characters \\
\texttt{memory\_before} & the stored record: \texttt{principal},
\texttt{memory\_slot}, \texttt{standing\_value} \\
\texttt{dialogue} & 9, 11, or 13 turns alternating USER/ASSISTANT, each
with a \texttt{turn\_id} \\
\texttt{evaluation\_query} & fixed instruction naming the slot and the
three permitted actions \\
\texttt{gold\_action} & KEEP, SUPERSEDE, or ABSTAIN \\
\end{longtable}
}

The closed schema is a design commitment, not housekeeping. There is no
conflict-type field, no design target, no component label, no annotator
note --- no channel of the kind that leaked the answer in v2. The
identifier is opaque so that it encodes nothing recoverable.

The dialogue has a fixed rhetorical shape: early turns state the
standing value and name the candidate value; a middle turn declares
which positions in a numbered evidence list are \textbf{indexed} to the
target slot and warns that the remaining positions belong to
neighbouring records; the evidence list is read out; a closing turn
restates the requested value. The decision must be made from the stored
record and this dialogue alone, before any content question is asked.

\subsection{5.9.4 The semantic component --- the load-bearing
idea}\label{the-semantic-component-the-load-bearing-idea}

A \textbf{semantic component} is a group of nine items sharing one
\texttt{principal}, one \texttt{memory\_slot}, one
\texttt{standing\_value}, and one candidate value. The nine form a 3 × 3
grid: three dialogue lengths (9, 11, 13 turns) crossed with the three
gold actions. Because the grid is crossed, both margins are exactly
balanced. Direct verification on the released public splits confirms
this:

\begin{itemize}
\tightlist
\item
  all \textbf{40/40} public components are exactly 3 KEEP / 3 SUPERSEDE
  / 3 ABSTAIN;
\item
  each dialogue length is exactly \textbf{40 KEEP / 40 SUPERSEDE / 40
  ABSTAIN}.
\end{itemize}

Neither verbosity nor length carries any signal about the action --- a
shortcut that trips many synthetic benchmarks is closed by construction
rather than by audit.

What varies inside a component is narrower still, and this is the
mechanism that makes the instrument work. All nine items draw their
evidence from a \textbf{byte-identical pool of six sentences} ---
verified for all \textbf{120/120} (component, length) triples in the
public data. The three actions differ \emph{only} in \textbf{which two
of those six sentences the dialogue indexes to the target slot}. Across
the 80 triples where the index statement parses mechanically, the three
indexed pairs are always distinct and always \textbf{partition the
six-sentence pool exactly} (80/80).

\subsection{5.9.5 A worked component}\label{a-worked-component}

For \texttt{principal\ =\ Rohan},
\texttt{memory\_slot\ =\ status\ message},
\texttt{standing\_value\ =\ available}, candidate value
\texttt{do\ not\ disturb}, all nine items contain the same six sentences
and the same closing turn. Only the index differs:

{\def\LTcaptype{none} % do not increment counter
\begin{longtable}[]{@{}
  >{\raggedright\arraybackslash}p{(\linewidth - 2\tabcolsep) * \real{0.5000}}
  >{\raggedright\arraybackslash}p{(\linewidth - 2\tabcolsep) * \real{0.5000}}@{}}
\toprule\noalign{}
\begin{minipage}[b]{\linewidth}\raggedright
Indexed pair (the two sentences bound to the slot)
\end{minipage} & \begin{minipage}[b]{\linewidth}\raggedright
Gold
\end{minipage} \\
\midrule\noalign{}
\endhead
\bottomrule\noalign{}
\endlastfoot
``The do-not-disturb period has an expiry record.'' + ``After the
expiry, Rohan writes that his status is available.'' & \textbf{KEEP} \\
``Rohan writes that the temporary period has ended.'' + ``He then
changes the status entry to do not disturb.'' & \textbf{SUPERSEDE} \\
``The retained records do not include the offsets needed to compare
them.'' + ``The expiry time and the status-change time use different
time zones.'' & \textbf{ABSTAIN} \\
\end{longtable}
}

The reasoning each pair demands differs in kind, which is the point:

\begin{itemize}
\tightlist
\item
  \textbf{KEEP} --- a temporary override expired and the standing value
  was reasserted \emph{after} it. The proposal is stale; the stored
  record is already correct.
\item
  \textbf{SUPERSEDE} --- the temporary period ended and a durable change
  followed. The proposal is current; the stored record is out of date.
\item
  \textbf{ABSTAIN} --- the two timestamps are in different time zones
  and the retained records lack the offsets needed to compare them.
  Neither conclusion is warranted. Note the character of this case: the
  evidence is not \emph{weak} or \emph{sparse}, it is
  \textbf{unresolvable}. This is deliberate, and it is why the 300
  ABSTAIN items carry a dual ambiguity flag: they encode designed
  unresolvability, not annotator uncertainty.
\end{itemize}

The other four sentences stay in every item as \textbf{controlled
distractors}. A system that reads the whole transcript rather than the
indexed evidence will find apparent support for all three actions in all
three items. The confound is not removed by paraphrasing; it is removed
by holding the surface form fixed and moving only the pointer. That is
the sense in which v4.3 is a \emph{controlled} synthetic benchmark
rather than merely a synthetic one.

\subsection{5.9.6 Splits and determinism}\label{splits-and-determinism}

Components, never items, are assigned to splits: train 180 items / 20
components, development 180 / 20, hidden 540 / 60. Each of the 20
domains contributes one component to train, one to development, and
three to hidden. Assignment uses recorded seed 4301. The build is
byte-deterministic --- archive timestamps are pinned to a constant
(1980-01-01) so identical inputs yield identical hashes. Because a
component never straddles a boundary, no system can observe the KEEP
variant of a scenario in training and the SUPERSEDE variant at test.

\subsection{5.9.7 The release gates are executable, not
narrative}\label{the-release-gates-are-executable-not-narrative}

\textbf{Shortcut audit.} Fitted on train, evaluated on hidden,
component-held-out, repeated over five seeds (4301--4305). It trains
word-level and character-level TF-IDF logistic regressions over
\textbf{twelve restricted views} of each item, plus punctuation-only and
length-only numeric features, plus decision trees at depths one through
five. Each view carries a pre-set accuracy ceiling:

{\def\LTcaptype{none} % do not increment counter
\begin{longtable}[]{@{}lrlr@{}}
\toprule\noalign{}
View & Ceiling & View & Ceiling \\
\midrule\noalign{}
\endhead
\bottomrule\noalign{}
\endlastfoot
\texttt{full\_text} & 0.65 & \texttt{assistant\_only} & 0.45 \\
\texttt{dialogue\_only} & 0.65 & \texttt{first\_three} & 0.50 \\
\texttt{target\_evidence\_only} & 0.65 & \texttt{query\_only} & 0.40 \\
\texttt{last\_two} & 0.50 & \texttt{memory\_only} & 0.40 \\
\texttt{candidate\_only} & 0.45 & \texttt{id\_only} & 0.40 \\
\texttt{previous\_turn} & 0.45 & \texttt{punctuation\_only} & 0.40 \\
\texttt{user\_only} & 0.45 & \texttt{length\_only} & 0.40 \\
\texttt{depth3\_tree} & 0.60 & & \\
\end{longtable}
}

The release ships with \texttt{failed\_gates\ =\ {[}{]}}: no probe
cleared its threshold. Two views deserve specific note.
\texttt{id\_only} confirms the opaque identifiers leak nothing --- the
exact class of failure that made v2's conflict-type field fatal.
\texttt{target\_evidence\_only} is the strictest self-criticism in the
build: it checks that even the \emph{correctly indexed} evidence,
stripped of dialogue context, does not linearly separate the three
classes, so the label depends on reading the evidence in context rather
than on a lexical signature of the evidence itself.

\textbf{Duplicate audit.} SHA-256 over seven restricted views, failing
on \emph{any} cross-split collision in the five decision-relevant ones
(\texttt{full\_text}, \texttt{dialogue\_only}, \texttt{candidate\_only},
\texttt{previous\_turn}, \texttt{last\_two}), plus character TF-IDF
(4--6 grams) cosine similarity over every cross-split pair, failing on
any edge at ≥ 0.995. Because components share an evidence pool by
construction, within-split near-duplication is expected and permitted;
the gate targets leakage across the split boundary specifically.

\textbf{Public/private boundary.} Public model-facing records exclude
gold labels, design targets, component metadata, private identifiers,
and annotator metadata. Hidden inputs, hidden gold, final human gold,
completed annotation packets, private mappings, and the detailed
restricted-view audits live only in a gitignored local boundary. Content
is addressed by two Merkle roots --- final gold \texttt{437b2a26…} and
annotations \texttt{a7e59f88…} --- and each distributed archive carries
its own SHA-256.

\subsection{5.9.8 Human validation and the weight it can
bear}\label{human-validation-and-the-weight-it-can-bear}

Two independent annotation packets cover all 900 items. The builder does
not accept them on trust. Every row must carry a schema version; a
\texttt{content\_hash} matching the coordinator's binding for that item;
an action in the permitted set; a \textbf{non-empty list of evidence
turn identifiers that all reference real turns in that dialogue}; a
rationale of at least 20 characters; a confidence in {[}0,1{]}; boolean
ambiguity and exclusion flags; an annotator identifier; and a
timezone-aware timestamp. Coverage must be exactly 900 with no duplicate
annotation identifiers. Any violation raises \texttt{RuntimeError} and
aborts the build.

The evidence-turn requirement is the substantive one: it makes an
annotation \emph{checkable against the item} rather than a bare label,
and it is what allows the builder to confirm that annotators engaged
with the indexed evidence structure rather than pattern-matching the
scenario.

Outcome: agreement on every action (raw 1.0, Cohen's κ 1.0), zero
exclusions, zero label adjudications, mean confidence 0.96 and 0.967.
Both annotators flagged all 300 ABSTAIN items as intentionally ambiguous
scenarios.

\textbf{What this does and does not establish.} Perfect agreement on a
task whose labelling rule is deterministic by construction certifies
that the manual is followable and that the items instantiate it. It does
\textbf{not} independently validate the construct: two annotators
applying the same unambiguous rule will agree whether or not the rule
measures memory governance. This is the honest reading, and it is why κ
= 1.0 appears in this thesis as a construction check rather than as
evidence of difficulty or of external validity. Completion is accepted
through dataset-owner attestation; the builder validates records,
coverage, hashes, and agreement, but does not authenticate reviewer
identity or annotation duration.

\subsection{5.9.9 Why the design makes errors
interpretable}\label{why-the-design-makes-errors-interpretable}

Because the three actions within a component are surface-matched, an
error is \emph{diagnostic} rather than merely wrong. Predicting
SUPERSEDE on the KEEP variant means the system read a sentence that was
present in the item but indexed to a different slot --- distractor
capture, specifically.

This is testable, and it was tested. Of the 28 errors the frozen
controller makes on the development split, 18 can be mechanically
resolved against a sibling item's indexed pair (same component, same
dialogue length, gold action equal to the \emph{predicted} action). In
\textbf{all 18}, the evidence supporting the predicted action was
present in the item but not indexed to the target slot. Every resolvable
error is distractor capture --- precisely the failure mode the component
design was built to isolate. The instrument is measuring what it was
constructed to measure, which is a stronger statement than any accuracy
figure in Chapter 7.

\subsection{5.9.10 What the construction cannot
deliver}\label{what-the-construction-cannot-deliver}

The control that makes v4.3 informative is bought by giving up realism,
and the trade must be stated plainly. Deliberately constructed scenarios
do not reproduce the distribution of natural user interactions; the
dialogues are stylized evidence-list recitations, not conversation. The
design is single-slot and single-step, so it cannot test interference
between records or multi-turn write trajectories. It is English-only.
Human completion is owner-attested rather than identity-authenticated.
And the release has no selected public redistribution license, so it
must not be publicly hosted until licensing is resolved. MGP-6 v4.3 is a
controlled diagnostic that must be paired with an external benchmark
such as LongMemEval (§7.11), never a substitute for one.

\section{6.1 Evidence policy: the ledger and claim
eligibility}\label{evidence-policy-the-ledger-and-claim-eligibility}

Run provenance comes from \texttt{artifacts/results\_ledger.json} v1.1;
scientific interpretation is superseded by
\texttt{artifacts/analysis/mgp6\_validity\_audit.json}. The failed
shortcut gate makes MGP-6 v2 policy rows diagnostic regardless of their
older ledger label. Mapping-dependent adapters, incomplete runs, and
every LongMemEval row remain non-primary. Annotation is described as one
independent annotator plus a separate model panel, never human IAA.

\section{6.2 Systems under test (MGP-6 v2, n =
180)}\label{systems-under-test-mgp-6-v2-n-180}

{\def\LTcaptype{none} % do not increment counter
\begin{longtable}[]{@{}
  >{\raggedright\arraybackslash}p{(\linewidth - 2\tabcolsep) * \real{0.5000}}
  >{\raggedright\arraybackslash}p{(\linewidth - 2\tabcolsep) * \real{0.5000}}@{}}
\toprule\noalign{}
\begin{minipage}[b]{\linewidth}\raggedright
System
\end{minipage} & \begin{minipage}[b]{\linewidth}\raggedright
Description
\end{minipage} \\
\midrule\noalign{}
\endhead
\bottomrule\noalign{}
\endlastfoot
B1 no-memory & answers without a store \\
B2 flat retrieval & stores everything, retrieves by similarity; no write
decisions \\
B3 recency & prefers most recent record; no abstention \\
B4 naive overwrite & new information always replaces old; no
abstention \\
Reference policy & deterministic lexical predicate policy used by the
released MGP-6 prediction artifact; not the prototype five-feature
scorecard \\
LLM controller & full ASK with qwen/qwen3.5-9b as the decision policy \\
Zero-shot LLM decision & same model asked for the action without the
governance scaffold \\
Few-shot (6-shot) LLM decision & 174/180 predictions; six missing IDs
scored INVALID for a coverage-aware bound (0.7833 / 0.7692);
\texttt{incomplete} \\
External adapters & Mem0 2.0.11, Letta (MemGPT), Graphiti 0.29.2, MemOS
via adapter harnesses \\
\end{longtable}
}

External-adapter scoring requires an action mapping because these
systems emit NOOP rather than ABSTAIN. Those mapped rankings are
retained in the artifact audit but excluded from primary comparisons. A
successor study must use each system's native observable operations and
resulting state.

\section{6.3 Ablations}\label{ablations}

Four single-component removals from the full rule controller, same items
and scaffold: \textbf{no\_abstain} (forced binary commit),
\textbf{no\_provenance}, \textbf{no\_typing}, \textbf{no\_lineage}.
Their within-generator failure signatures (§7.3) are mechanism
diagnostics: they show how this implementation changes when a switch is
removed, but do not independently establish that each component is
necessary in other datasets or controllers.

\section{6.4 LongMemEval-M protocol}\label{longmemeval-m-protocol}

Stratified 180-item slice (6 × 30, seed 0). Five systems (B1--B4,
controller) were run under the v7 configuration (answerer qwen3.5-9b,
judge google/gemma-4-31b-it); the run is \textbf{diagnostic only} due to
documented provider contamination (§7.6). The v8 failure replay
re-executed exactly the 70 v7 controller failures under a changed
configuration and is \texttt{failure\_replay\_only}, not an evaluation
of the full slice. The complementary Modal top-up targeted 110 prior v7
successes but completed only 84, so it is \texttt{incomplete}; it
nevertheless provides decisive counter-evidence to carrying all
successes forward (§7.7). The frozen final-1.1 protocol specifies a
homogeneous full-180 rerun with qwen/qwen3.5-9b and a Gemma judge under
a fixed provider and retry policy. Its API smoke failed before inference
because the supplied OpenRouter credential returned HTTP 403, so the run
is \texttt{externally\_pending} and no LongMemEval value is final.

\section{6.5 Metrics}\label{metrics}

Decision layer: action accuracy, macro-F1, per-action F1, confusion
counts. Commit safety: invalid-commit rate (share of gold-ABSTAIN items
on which the system commits a write). Memory state: lineage integrity;
stale-retrieval rate. Probe behavior: premise resistance; implicit
application. Downstream (LME-M): judge-graded answer accuracy with
per-competency decomposition and stratified bootstrap CIs. Formal
definitions in Appendix C. Excluded as uninformative: the shipped
over-/under-forgetting rates (§5.8). Premise resistance and
invalid-commit are bespoke metrics defined in \texttt{metrics\_mgp6.py}
and are interpreted descriptively, not as externally validated
constructs (threat C5).

\section{6.6 Statistical analysis}\label{statistical-analysis}

The primary analysis uses the held-out 90-item test split and
10,000-replicate paired cluster bootstrap over
\texttt{template\_family\_id}, repeated with minimal-pair groups plus
singleton items. Both item-weighted and family-macro estimates are
reported. Exact item-level McNemar is secondary because within-family
independence is doubtful. Zero event counts receive one-sided
finite-sample upper bounds. No superiority, equivalence, or
non-inferiority language is licensed without a pre-specified contrast
and margin. On LongMemEval, no hypothesis test is computed because no
clean comparative design exists. Reproduction:
\texttt{python3\ scripts/analysis/audit\_mgp6\_validity.py}.

\section{6.7 Leakage and researcher-degrees-of-freedom
controls}\label{leakage-and-researcher-degrees-of-freedom-controls}

Train+development lookup probes over type, conflict type, and public
metadata precede policy interpretation. Conflict type reaches 1.000 test
accuracy and therefore fails the construct-validity gate: test-first
split hygiene alone is insufficient. The v7→v8 fix loop had no
confirmation set, so its output is \texttt{failure\_replay\_only};
overrides remain separate from raw results; and invalid, incomplete,
mapped, and superseded runs are excluded from primary claims.

\section{6.8 MGP-6 v4.3 experimental protocol and completed hidden
run}\label{mgp-6-v4.3-experimental-protocol-and-completed-hidden-run}

MGP-6 v4.3 uses train for loader and engineering checks, and development
for prompt, model, configuration, ablation, cost, and latency selection.
Each system submits one strict record,
\texttt{\{"public\_id":"MGP43-...","action":"KEEP"\}}; the validator
rejects missing, duplicate, unknown, extra, and invalid-action records.
Public loading and prompt construction are deterministic, prediction
writing is resume-safe, and raw-response logs record latency, token,
retry, and failure accounting without private fields.

Before hidden prediction, the release and split hashes, provider and
exact model ID, prompt and system-prompt hashes, decoding, seed, retry
policy, timeout, concurrency, evaluator version, code commit,
environment lock, and prediction schema were frozen and recorded. The
selected public-indexed-evidence ASK live controller with
qwen/qwen3.5-9b was then run exactly once on the 540-item hidden split.
Its prediction file was complete, duplicate-free, and sealed with
SHA-256 e6a598799115c9fc80ae3e731f1b5d59a4947e48df5356b12ba518bafb846806
before the local-private scorer opened hidden gold. A transport
disconnect after 432 durable predictions was recorded; the frozen run
resumed only the remaining 108 records and did not change controls or
duplicate IDs.

The private evaluator reports accuracy, macro-F1, per-class
precision/recall/F1, confusion matrix, coverage, selective and
abstention metrics, and a 5,000-replicate component-cluster bootstrap
with seed 4301. It does not expose gold records in its aggregate output.
Provider telemetry supplies latency, tokens, failures, retries, and
invalid-response counts, but no price field; cost is therefore recorded
as unavailable. LongMemEval remains a separate external evaluation, not
a substitute for this controlled diagnostic. Because only one final v4.3
configuration was pre-registered, there is no final hidden baseline,
ablation, or paired comparison.

\subsection{6.8.1 Systems, metrics, and statistical treatment for
v4.3}\label{systems-metrics-and-statistical-treatment-for-v4.3}

The v4.3 comparison covers twelve configurations on the public
development split: the pre-registered frozen controller; the successor
controller in both its as-run and guard-composed forms; four structural
baselines (B1 no memory, B2 flat retrieval, B3 recency, B4 naive
overwrite); an unstructured zero-shot decision policy; the deterministic
public-evidence rule policy in isolation; and four component ablations
(no abstention, no provenance, no typing, no lineage). Three deployed
external memory systems --- Letta, Mem0, and Graphiti --- were run
through adapters written for this thesis. Because their native
operations do not match the KEEP/SUPERSEDE/ABSTAIN vocabulary, their
scores depend on an action mapping and are reported as harness
diagnostics, never as capability measurements of those systems in their
intended use. All model-backed rows use \texttt{qwen/qwen3.5-9b} at
temperature 0, seed 4301, 128 maximum new tokens, and concurrency 1.

Primary decision metrics are action accuracy, macro-F1, per-action
precision/recall/F1, Cohen's κ, and the multiclass Matthews correlation.
The latter two are interpretable here precisely because every split is
exactly action-balanced, so neither is inflated by a class prior.

Errors are decomposed by \emph{operational harm} rather than reported as
a single rate. Four quantities are tracked, each with its own
denominator: action false commit on gold-ABSTAIN items, separated into
the state-preserving ABSTAIN→KEEP and state-corrupting ABSTAIN→SUPERSEDE
components; missed update on gold-SUPERSEDE items; destructive overwrite
of a gold-KEEP record; and unnecessary abstention over items that
required a decision. Every rate carries a Wilson 95\% interval, and a
zero count is reported with a finite-sample upper bound rather than as
zero risk. Expected loss is computed under balanced, commit-averse, and
update-responsive cost matrices fixed before scoring.

Hidden-split uncertainty uses a 5,000-replicate component-cluster
bootstrap at seed 4301, because items within a semantic component share
a scenario skeleton and are not independent. On the public development
split, component-cluster resampling is \textbf{unavailable} --- the
component map is private and the public scorer does not read it --- so
those intervals are 10,000-replicate item bootstraps and are
consequently \emph{narrower than a component-aware estimate would be}.
Paired system contrasts on development use exact two-sided McNemar tests
with Holm correction over the reported family; they inherit the same
independence assumption. This thesis relies on them only for contrasts
whose effect sizes (b/c splits of 98/6 and larger) could not plausibly
be produced by within-component dependence, and flags the one marginal
contrast --- the 0.117 successor gain --- as the place where the
assumption matters most.

Operational accounting is taken from per-response records rather than
harness wall-clock timing: provider-reported latency per response (mean,
median, p95, maximum), prompt and completion tokens, transport failures,
in-client retries, and invalid-response counts. Monetary cost is
reported only where the provider emitted a price field; on every run in
this thesis it did not, and cost is recorded as unavailable rather than
reconstructed from list-price arithmetic.
